\documentclass[10pt,a4paper]{amsart}
\usepackage[margin=19mm]{geometry}
\usepackage{amsmath,amssymb,mathtools}
\usepackage[scr=rsfs,cal=euler]{mathalfa}
\usepackage{xfrac}
\usepackage[unicode,hidelinks,bookmarks=false]{hyperref}
\newcommand{\ie}{i.e.\ }
\newcommand{\cf}{cf.\ }
\newcommand{\resp}{respectively\ }
\newcommand{\Subsection}{\subsection}
\providecommand{\nobreakdash}{\nobreak\hbox{-}}
\setlength{\emergencystretch}{2em}
\multlinegap0pt
\usepackage{bm}
\usepackage{bbm}
\usepackage{dsfont}
\usepackage{xspace}
\usepackage{cancel}
\usepackage{mathtools}
\usepackage{amsthm}
\makeatletter
\@ifundefined{thm}{\newtheorem{thm}{Thm}}{}
\@ifundefined{lem}{\newtheorem{lem}[thm]{Lemma}}{}
\makeatother
\DeclareMathOperator{\mult}{mult}
\DeclarePairedDelimiter{\abs}{\lvert}{\rvert}
\DeclarePairedDelimiter{\floor}{\lfloor}{\rfloor}
\newcommand{\cA}{\mathcal{A}}
\newcommand{\cG}{\mathcal{G}}
\newcommand{\cS}{\mathcal{S}}
\newcommand{\multibinom}[2]{\left(\kern-0.3em\left(\genfrac{}{}{0pt}{}{#1}{#2}\right)\kern-0.3em\right)}
\title[On the Thickness of Infinite Generalized Sidon Sets, II]{On the Thickness of Infinite Generalized Sidon Sets, II}
\author{Kevin O'Bryant}
\date{Source: arXiv:2607.23795v1 (CC BY 4.0)}
\begin{document}
\maketitle

\noindent\emph{Source and licence.} On the Thickness of Infinite Generalized Sidon Sets, II. Version arXiv:2607.23795v1 (CC BY 4.0), distributed under the Creative Commons Attribution 4.0 International licence, as recorded in the arXiv licence metadata for this exact version and shown on the abstract page https://arxiv.org/abs/2607.23795v1. Full rights notice: licenses/arxiv-2607.23795-CC-BY-4.0.txt.\par\smallskip

\noindent\textit{Editorial scope.} Counted unit: the Lem whose source label is \texttt{lem:sumset energy bound} in the article (source0.tex lines 530--540, source label \texttt{lem:sumset energy bound}), delivered together with the article's own complete original proof of it (source0.tex lines 542--642, one \texttt{proof} environment of 101 lines). No mathematics is written, repaired or re-derived: statement and proof are the article's own text line for line. Required context retained uncounted: lem:cancel (lines 195--197); lem:B\_h is B\_p (lines 200--202); lem:easy Bh bound (lines 254--256); eq:factor Phi (lines 462--464); lem:Jia in action (lines 524--528). Every definition, display and label of those lines is retained unchanged. Omitted from this package and not redistributed: 1--194, 198--199, 203--253, 257--461, 465--523, 529--529, 541--541, 643--851. Nothing inside a retained excerpt is omitted or altered: every retained body is delivered exactly as the article's lines are written, so no omission receipt of any kind accompanies this package. Citations: the retained excerpts carry no citation command at all, so no bibliography entry of the article is transcribed and no reference is redistributed. Reference adaptation: none was needed and none was made; every reference inside the delivered unit resolves inside it, so no unresolved reference or citation is produced. Printed slips kept literal: no printed word, symbol, hypothesis or number of the counted statement or of its proof was altered. Layout and class: the article's own class is \texttt{article} with packages including amsmath, amssymb, amsthm, mathtools, enumitem, geometry, hyperref, amsrefs, tikz, pgfplots, todonotes; the wrapper is the lane's pinned \texttt{amsart} editorial wrapper at 10pt on A4, which restates the theorem-like environments the retained text needs and adds the article's own macro definitions that the retained text needs (\texttt{\textbackslash abs}, \texttt{\textbackslash cA}, \texttt{\textbackslash cG}, \texttt{\textbackslash cS}, \texttt{\textbackslash floor}, \texttt{\textbackslash mult}, \texttt{\textbackslash multibinom}), with extra wrapper packages (bm, bbm, dsfont, xspace, cancel, mathtools). The delivered numbering is the unit's own. Assets: no retained span contains a figure environment, a graphics-inclusion command, a PDF-page inclusion, a table or a TikZ picture; no image file is copied into this package, the package declares no asset dependency and no image file is redistributed with it. Licence: version arXiv:2607.23795v1 of this article is distributed under the Creative Commons Attribution 4.0 International licence, as recorded in the arXiv licence metadata for this exact version and shown on the abstract page https://arxiv.org/abs/2607.23795v1. A full rights notice is delivered at licenses/arxiv-2607.23795-CC-BY-4.0.txt. Characters: every retained excerpt is pure ASCII, so the delivered engine, XeLaTeX with its default Latin Modern text font, reports no missing character.
\begin{lem}\label{lem:cancel}
Let $P,P',Q,Q'$ be multisets with $P\cap Q=P'\cap Q'=\emptyset$.  If $P\uplus Q'=P'\uplus Q$, then $P=P'$ and $Q=Q'$.
\end{lem}

\begin{lem}\label{lem:B_h is B_p}
  Suppose that $\cA$ is a $B_h$-set and $1\le p \le h$. Then $\cA$ is a $B_p$-set. 
\end{lem}

\begin{lem}\label{lem:easy Bh bound}
  If $\cG$ is a $B_h$-set and $m\ge 1$, then $G(m) \le (h\cdot h!)^{1/h} \cdot m^{1/h}$.
\end{lem}

\begin{equation}\label{eq:factor Phi}
   \abs*{\Phi(r ,p,q;x,L)} = \abs*{\cA_L}^r \abs*{T(p,q;x,L)}.
\end{equation}

\begin{lem}\label{lem:Jia in action}
For all $N\ge 1$ and $1\le r <k$, one has
  \[\abs*{\Phi(2r,k-r,k-r;0,N)} \le N \cdot C(2r,k-r,k-r).\]
  %(2r +1) \frac{(k+r)!}{(k-r)!} \sum_{j=0}^{k-r}(4r)^{j}.\]
\end{lem}

\begin{lem}\label{lem:sumset energy bound}
Let $h$ be even, let $\cA$ be a $B_h$-set satisfying 
  \begin{equation}\label{eq:tau appears}
    A(n) \ge \tau \, \left(\frac{n}{\log n}\right)^{1/h}
  \end{equation}
for some $\tau\ge 0$ and all $n\ge n_0\ge 3$. let $k \coloneqq h/2$, and let $\cS=k\cA$. For each $N\ge n_0$, there is a $t^\ast\in [0,N)$ that makes the block counts 
  \[F_\ell\coloneqq S(t^\ast+\ell N) - S(t^\ast+(\ell-1)N)\]
    %\abs*{\cS\cap[t^\ast+(\ell-1)N, t^\ast+\ell N)}\] 
satisfy, with $M\coloneqq \floor{N/\log^3(N)}$,
  \[   \sum_{\ell=1}^{M}\binom{F_\ell}{2} \le \frac{1}{2}\,N + o(N).\]
\end{lem}

\begin{proof}
For $N\ge 3$, put $W\coloneqq (M+1)N=N^2/\log^3(N)+O(N)$.  For integers
$t\in[0,N)$ and $\ell \in [1,M]$, set
  \[F_\ell^{(t)}\coloneqq \abs*{\cS \cap [t+(\ell-1)N, t+\ell N)}
  =S(t+\ell N)-S(t+(\ell-1)N).\]  
A pair $(s,s') \in \cS_W\times \cS_W$ with $s-s'=d \in[1, N)$, lies in a common block $[t+(\ell-1)N, t+\ell N)$ for at most $N-d$ offsets, so
\begin{equation}\label{eq:average}
   \frac1N\sum_{t=0}^{N-1}\sum_{\ell=1}^{M}\binom{F_\ell^{(t)}}{2}
    \le \frac1N\sum_{d=1}^{N-1}(N-d) P(d),
\end{equation}
where $P(d)$ is the number of such pairs with difference $d$.  
That is,
  \[P(d) \coloneqq \#\left\{ (s,s') \in (\cS_W)^2 \colon s-s'=d \right\}.\]
Choose $t^\ast\in[0,N)$ so that
  \[\sum_{\ell=1}^{M}\binom{F_\ell^{(t^\ast)}}{2} \le \frac1N\sum_{t=0}^{N-1}\sum_{\ell=1}^{M}\binom{F_\ell^{(t)}}{2}\]
and set $F_\ell \coloneqq F_\ell^{(t^\ast)}$. We have
  \begin{equation}\label{eq:P appears}
    \sum_{\ell=1}^{M}\binom{F_\ell}{2}
    \le \frac1N\sum_{d=1}^{N-1}(N-d) P(d).
  \end{equation}

For each $s\in \cS$, we identify the unique $V_s \in \multibinom{\cA}{k}$ with $\Sigma V_s =s $ (unique because $\cA$ is a $B_k$-set by Lemma~\ref{lem:B_h is B_p}). 
We stratify $P(d)$ by the size of $V_s \Cap V_{s'}$. Namely, set
  \[P_r(d) \coloneqq \#\left\{ (s,s')\in (\cS_W)^2 \colon s-s'=d ,\abs*{V_s \Cap V_{s'}}=r \right\}.\]
We have, $P(d) = \sum_{r=0}^k P_r(d)$. For $d\ge1$, we have $P_k(d)=0$, as $|V_s \Cap V_{s'}| =k$ implies that $V_s=V_{s'}$, so that $s=s'$ and $s-s'=d=0$. 
Thus, Line~\eqref{eq:P appears} becomes
  \begin{equation}\label{eq:P_r appears}
    \sum_{\ell=1}^{M-1}\binom{F_\ell}{2}
    \le \frac1N\sum_{d=1}^{N-1}(N-d) P_0(d)+ \sum_{r=1}^{k-1} \frac1N\sum_{d=1}^{N-1}(N-d) P_r(d).
  \end{equation}

We now show that for $d\ge 1$, we have $P_0(d)\le 1$. 
We see that $P_0(d)$ counts the number of pairs $(s,s')\in (\cS_W)^2$ with $s-s'=d$ and $V_s \Cap V_{s'}=\emptyset$.
Suppose that both $(s,s')$ and $(u,u')$ are such pairs. 
Then $s+u'=u+s'$, an identity of $2k$-fold sums, so by the $B_{2k}$ property
  \[V_s \uplus V_{u'} = V_{u} \uplus V_{s'}.\]
By Lemma~\ref{lem:cancel}, $V_s=V_u$ and $V_{s'}=V_{u'}$, from which it follows that $s=u$ and $s'=u'$. 

Therefore
\begin{equation}\label{eq:disjoint}
   \frac1N \sum_{d=1}^{N-1}(N-d) P_0(d) \le \frac 1N \sum_{d=1}^{N-1}(N-d) = \frac12 (N-1).
\end{equation}
The bound~\eqref{eq:P_r appears}, with $N-d \le N$ becomes
  \begin{equation}\label{eq:P_0 is gone}
    \sum_{\ell=1}^{M-1}\binom{F_\ell}{2}
    \le \frac12 N + \sum_{r=1}^{k-1} \sum_{d=1}^{N-1} P_r(d).
  \end{equation}

We now consider $1 \le r < k$, and show that for each such $r$ we have $\sum_d P_r(d)=o(N)$.  
Consider a pair $(s,s')$ counted by $\sum_{d=1}^{N-1} P_r(d)$, so that there are unique $\alpha\in \multibinom{\cA}{r}, \beta,\delta \in \multibinom{\cA}{k-r}$ with $\alpha =V_s \Cap V_{s'}$ (so $\abs{\alpha}=r$), $V_s = \alpha \uplus \beta$ and $V_{s'} = \alpha \uplus \delta$ (so $\abs{\beta}=k-r = \abs{\delta}$). 
Note that $\beta \cap \delta = \emptyset$, since at each $x$ either $\mult_\beta(x)=0$ or $\mult_\delta(x)=0$.
The map from $(s,s')$ to $(\alpha;\beta,\delta)$ is injective. 
Since $\Sigma \beta -\Sigma\delta =s-s' =d \in [1,N)$, we see that $(\beta,\delta) \in T(k-r,k-r;0,N)$.

Hence, using $r\ge 1$ we have the bound $\binom{A(W)+r-1}{r}\le A(W)^{r}$, and so
\[
   \sum_{d=1}^{N-1}P_r(d)
   \le \abs*{\multibinom{\cA_W}{r}}\cdot\abs*{T(k-r,k-r;0,N)}
   \le A(W)^{r} \abs*{T(k-r,k-r;0,N)}.
\]
Now multiply and divide by $A(N)^{2r}$ and apply the factorization~\eqref{eq:factor Phi} to get
\begin{align*}
   \sum_{d=1}^{N-1}P_r(d)
   &\le\frac{A(W)^{r}}{A(N)^{2r}}\cdot A(N)^{2r} \abs*{T(k-r,k-r;0,N)}\\
   &=\frac{A(W)^{r}}{A(N)^{2r}} \abs*{\Phi(2r,k-r,k-r;0,N)}.
\end{align*}
Lemma~\ref{lem:Jia in action} gave us
\[
   \abs*{\Phi(2r,k-r,k-r;0,N)}\le N \cdot C(2r,k-r,k-r),
\]
bringing our bound to
  \[\sum_{d=1}^{N-1} P_r(d) \le  C(2r,k-r,k-r) \, \frac{A(W)^{r}}{A(N)^{2r}}\,  N.\]


It remains to see that $A(W)^r/A(N)^{2r}=o(1)$. This is why Jia introduced a growth regularity hypothesis $A(N^2) = O(A(N)^2)$. As Erd\H{o}s, Helm, and Jia weren't thinking of producing explicit constants, they never separated $M$ to be its own parameter, and just used $M=N$. The flexibility afforded by this additional parameter is what allows our argument to proceed without Jia's regularity hypothesis. That is, this step is why we have $M=N/\log^3 N$ and not $M=N/\log N$ (as in Part I of this series of papers) or $M=N$ (as in Chen's work).

By Lemma~\ref{lem:easy Bh bound} and
$W=N^2/\log^3 N$,
\[
   A(W)\le(h\cdot h!)^{1/h}W^{1/h}
   =(h\cdot h!)^{1/h} \frac{N^{2/h}}{\log^{3/h} N},
\]
while~\eqref{eq:tau appears} gives $A(N)\ge \tau N^{1/h}/\log^{1/h} N$ for
$N\ge n_0$.  Therefore
\[
   \frac{A(W)^{r}}{A(N)^{2r}}
   \le\frac{(h\cdot h!)^{r/h}}{\tau^{2r}} 
     \frac{\log^{-3r/h} N}{\log^{-2r/h} N}
   =\left(\frac{(h\cdot h!)^{1/h} / \tau^{2}}{\log^{1/h} N}\right)^r,
\]
which goes to $0$ as $N\to \infty$.
Combining, for each fixed $1\le r\le k-1$,
  \[
   \sum_{d=1}^{N-1}P_r(d)
   \le\frac{(h\cdot h!)^{r/h} C(2r,k-r,k-r)}{\tau^{2r}}\cdot
     \frac{N}{\log^{r/h} N}=O\left(\frac{N}{\log^{1/h} N}\right),
  \]
so the finite sum $\sum_{r=1}^{k-1}\sum_{d}P_r(d)=o(N)$.  The bound on Line~\eqref{eq:P_0 is gone} becomes
  \[   \sum_{\ell=1}^{M-1}\binom{F_\ell}{2}\le\frac{N}{2}+o(N), \]
which concludes the proof of Lemma~\ref{lem:sumset energy bound}.
\end{proof}

\end{document}
